\section{Conclusion}
\label{sec:conclusion}

HyNoC is a hybrid circuit-switch/wormhole Network-on-Chip designed for distributed
VLIW computing on FPGA. By combining source routing---which establishes a dedicated
path through a distributed request/grant handshake---with wormhole flit streaming,
HyNoC achieves bounded, deterministic transfer latency without the area overhead of
virtual channels.

Congestion is managed statically by the compiler or system software, which selects
routes across the rich set of shortest paths available in mesh and
higher-dimensional topologies. The combinatorial analysis presented in
Section~\ref{sec:switching} shows that enriching the physical topology is more
area-efficient than adding virtual channels for the targeted workload class,
provided that a compile-time scheduler distributes traffic across available paths.

The parallel round-robin arbiter provides starvation-free, fixed-latency
arbitration at each egress port, and per-port independent clock domains allow the
NoC to interface cleanly with heterogeneous processing elements running at
different frequencies. The multiple local interface mechanism offers a practical
alternative to virtual channels when a node requires simultaneous non-blocking
communication with multiple peers.

Verilator co-simulation on a $4\times4$ mesh validates the analytical latency
model: measured return-path latency satisfies $L=12H+P-1$ (dual-clock) and
$L=9H+P-1$ (single-clock), with the $3H$ gap between modes explained exactly
by the CDC pipeline overhead eliminated in single-clock operation.  A
LLaMA~3~8B FFN up-projection benchmark (Q8\_0 weights, BF16 activations,
$14336\times4096$) with a single master completes in 43\,M cycles; a
four-master quadrant configuration, where each corner master owns a $2\times2$
traffic-isolated quadrant, reduces this to 8.6\,M cycles---a $5\times$ speedup
with zero cross-quadrant link utilization, confirming the deterministic
throughput scaling predicted by the link-disjoint topology design.

HyNoC is released as open-source hardware under the CERN-OHL-P~v2 license,
providing a complete RTL implementation with simulation testbenches and FPGA board
support.

Future work includes: formal deadlock-freedom analysis for arbitrary source route
sets; synthesis results and latency characterization on representative FPGA
devices (Section~\ref{sec:evaluation} provides Verilator-based characterization;
on-device measurement remains open); integration with a VLIW compiler back-end
for automatic route assignment and injection scheduling; and extension of the XY
routing protocol placeholder to a full hardware implementation.
